\newpage
\section{Proofs}\label{sec:tech:proof}

\subsection{Lemmas for s-LTL$_f$}

We introduce several auxiliary lemmas for s-LTL$_f$ formulae.

\begin{lemma}\label{lemma:allTL}[All well-fromed traces] $\denot{\allTL} = \{ tr ~|~ \wellfoundedvdash tr  \}$.
\end{lemma}

\begin{lemma}\label{lemma:espTL}[Epsilon set of trace] $\denot{\epsTL} = \{ [] \}$.
\end{lemma}

\begin{lemma}\label{lemma:empTL}[Empty set of trace] $\denot{\empTL} = \emptyset$.
\end{lemma}

\begin{definition}[Well-formed handler context]
  \label{lemma:built-in-typing-ap}
  The handler specification $\Delta$ and capability context $\Sigma$ is well-formed iff for all operator $\eff{op}$ and its \PAT $\overline{y{:}b} \garr \overline{x{:}t} \sarr \rg{H}{\lastA\msgB{op}{\overline{y}}{\phi}}{F}$ and capability $\Theta$ in $\Delta$ satisfying
{\small\begin{align*}
    \forall \overline{y{:}b}.\forall \alpha \in \denot{A}. \forall \overline{c \in \denotation{t}}. \forall \overline{\alpha_i}. \forall \overline{m_i}.\; \alpha_1 \listconcat [m_1] \listconcat ... [m_{n}] \listconcat \alpha_{n+1} \in \denot{F} \impl \alpha \vDash \zevent{op}{\overline{c}} \Downarrow \{m_i\}
\end{align*}}\noindent
also, $\overline{y{:}\rawnuot{b}{\top}},\overline{x{:}t} \vdash F \subseteq \bigwedge_{\eff{op} \in \Theta} \finalA \msg{op}{\top}$ should hold.
\end{definition}

\begin{theorem}\label{theorem:fundamental-ap}[Fundamental Theorem] Given
  a well-formed handler specification $\Delta$,
    the trace of effects produced by a well-typed term
    $e$ is captured by its corresponding \PAT $\tau$:
  $\Gamma; \Delta; \Theta \vdash e : \tau \implies \forall \sigma, \sigma \in
    \denotation{\Gamma} \impl \sigma(e) \in
    \denotation{\sigma(\tau)}$.\footnote{The proofs of theorems in this section are provided in the \techreport{}.}
\end{theorem}
\begin{proof} We proceed by induction over our type judgment $\Gamma;\Delta;\Theta \vdash e : \tau$, which has seven cases proved as following:
\begin{enumerate}[label=Case]
\item:
{\footnotesize
\mprooftr{35mm}
{}
{TRet}
{$\Gamma; \Delta; \emptyset \vdash () : \rg{H}{\epsTL}{F}$}
\ \\ \ \\}
This rule assume that $\Theta \equiv \emptyset, e \equiv (), A \equiv \epsTL$, thus we need to prove
{\footnotesize\begin{align*}
    \forall \sigma, \sigma \in \denot{\Gamma} \impl \sigma(()) \in \denot{\sigma(\rg{H}{\epsTL}{F})}
\end{align*}}\noindent
From the induction hypothesis and the precondition of this rule, we have
{\footnotesize\begin{alignat}{2}
    \setcounter{equation}{0}
        &\forall \sigma. \sigma(()) = () \\
        &\forall \sigma. \sigma(\rg{H}{\epsTL}{F}) = \rg{\sigma(H)}{\epsTL}{\sigma(F)}
\end{alignat}}\noindent
where we just need to prove the term $()$ is in the denotation of a \PAT in from $\rg{H}{\epsTL}{F}$, which means for all $\alpha_h\;\alpha\;\alpha_f\;\beta\;\beta'\;e_h\;e_f$,
{\footnotesize\begin{alignat}{2}
&\eraserf{\Gamma} \basicvdash () : \Code{unit} \ptag{from basic typing ??}
\\&\msteptr{[]}{(\emptyset, e_h)}{\alpha_h}{(\beta, ())} 
\land \msteptr{\alpha_h}{(\beta, ())}{\alpha}{(\beta', ())}
\land \msteptr{\alpha_h \listconcat \alpha }{(\beta', e_f)}{\alpha_f}{(\emptyset, ())}
\ptag{assumption of $\denot{\tau}$}
\\&\alpha_h \in \denot{H} \land \alpha_f \in \denot{F} \ptag{assumption of $\denot{\tau}$}
\end{alignat}}\noindent
Since there is no small-step reduction rule for the term $()$, thus the relation $\msteptr{\alpha_h}{(\beta, ())}{\alpha}{(\beta', ())}$ is derived from reflexivity case of multi-step reduction. Thus, $\alpha$ is empty trace $[]$, which included by the denotation of $\epsTL$ (\autoref{lemma:espTL}). Then the proof immediate holds in this case.

\item:
{\footnotesize
\mprooftr{73mm}
{
$\Gamma \vdash \eff{op} : \overline{x_i{:}t_i}\sarr \rg{H}{\lastA\msg{op}{\phi}}{A \seqA F}$ \quad
$\forall i. \Gamma \vdash v_i : t_i$ \quad
$\GEN(\eff{op})$ \quad $(\eff{op}, B') \in \Sigma$ \quad
 $\Gamma ~|~ B \cup B' \vdash e : \rg{H \seqA \lastA\msg{op}{\phi\msubst{x_i}{v_i}}}{A}{F}$
 }
{TGen}
{$\Gamma ~|~ B \vdash \zgen{op}{\overline{v_i}}{e} : \rg{H}{\lastA\msg{op}{\phi\msubst{x_i}{v_i}} \seqA A}{F}$}
\ \\ \ \\}
This rule assume that $e \equiv \zgen{op}{\overline{v}}{e}, \tau \equiv \rg{H}{\lastA\msg{op}{\phi\msubst{x}{v}} \seqA A}{F}$, thus we need to prove
{\footnotesize\begin{align*}
    \forall \sigma, \sigma \in \denot{\Gamma} \impl \sigma(\zgen{op}{\overline{v}}{e}) \in \denot{\sigma(\rg{H}{\lastA\msg{op}{\phi\msubst{x}{v}} \seqA A}{F})}
\end{align*}}\noindent
From the induction hypothesis and the precondition of this rule, we have
{\footnotesize\begin{alignat}{2}
\setcounter{equation}{0}
&\Gamma \vdash \eff{op} : \overline{x{:}t}\sarr \rg{H}{\lastA\msg{op}{\phi}}{A \seqA F} \ptag{assumption} \\
&\forall i. \Gamma \vdash v_i : t_i \ptag{assumption} \label{loc:1-1} \\
&\GEN(\eff{op}) \land (\eff{op}, B') \in \Sigma \ptag{assumption} \\
&\Gamma ~|~ B \cup B' \vdash e : \rg{H \seqA \lastA\msg{op}{\phi\msubst{x}{v}}}{A}{F} \ptag{assumption} \\
&\forall \sigma \in \denot{\Gamma}. \sigma(e) \in \denot{\sigma(\rg{H \seqA \lastA\msg{op}{\phi\msubst{x}{v}}}{A}{F})} \ptag{induction hypothesis} \\
&\forall i. \forall \sigma \in \denot{\Gamma}. \sigma(v_i) \in  \sigma(\denot{t_i}) \ptag{\autoref{loc:1-1} and ??} \\
&\msteptr{\alpha_h}{(\beta, \zgen{op}{\overline{v}}{e})}{\alpha}{(\beta', e)}  \ptag{??} \\
% &\forall \alpha \in \denot{A}. \forall \overline{c \in \denot{t}}. \forall \overline{\alpha_i}. \forall \overline{m_i}.\; \alpha_1 \listconcat [m_1] \listconcat ... [m_{n}] \listconcat \alpha_{n+1} \in \denot{F} \impl \alpha \vDash \zevent{op}{\overline{c}} \Downarrow \{m_i\} \ptag{??} \\
\end{alignat}}\noindent
\end{enumerate}
\end{proof}
